\documentclass[11pt,a4paper]{article}
\usepackage[utf8]{inputenc}
\usepackage{amsmath,amssymb,amsthm}
\usepackage{listings}
\usepackage[margin=2.5cm]{geometry}
\usepackage{xcolor}

\lstset{
  language=C++,
  basicstyle=\ttfamily\small,
  keywordstyle=\color{blue}\bfseries,
  commentstyle=\color{green!50!black},
  stringstyle=\color{red!70!black},
  numbers=left,
  numberstyle=\tiny\color{gray},
  breaklines=true,
  frame=single,
  tabsize=4,
  showstringspaces=false
}

\title{IOI 1992 -- Problem 2: Teletext}
\author{Solution}
\date{}

\begin{document}
\maketitle

\section{Problem Statement}

A teletext display renders characters using a fixed-size pixel grid (a
dot-matrix font). Given a library of known character patterns and an input
image, recognize each character cell and decode the displayed message.

Each character occupies a cell of $w \times h$ pixels. The input image is a
pixel grid whose dimensions are exact multiples of $w$ and $h$. Match each
cell against the known patterns to decode the text.

\section{Solution}

\subsection{Algorithm}

\begin{enumerate}
  \item \textbf{Load patterns:} For each known character, store its
    $h \times w$ pixel grid.
  \item \textbf{Segment the image:} Divide the input image into cells of
    size $h \times w$.
  \item \textbf{Match each cell:} Compare pixel-by-pixel against all
    patterns. For exact matching, find the pattern with zero differences.
    For noisy input, find the pattern with the smallest Hamming distance
    (number of differing pixels).
\end{enumerate}

\subsection{Hamming Distance}

The Hamming distance between two $h \times w$ binary grids $A$ and $B$ is:
\[
  d(A, B) = \sum_{r=0}^{h-1} \sum_{c=0}^{w-1}
  \mathbf{1}[A[r][c] \ne B[r][c]].
\]

For noise-free input, the correct character has $d = 0$.

\section{Complexity Analysis}

\begin{itemize}
  \item \textbf{Time:} $O(T \cdot K \cdot wh)$ where $T$ is the number of
    character cells, $K$ is the alphabet size, and $wh$ is the cell size.
  \item \textbf{Space:} $O(Kwh + WH)$ for storing patterns and the input
    image.
\end{itemize}

For typical parameters ($K \le 128$, $w, h \le 16$), this is efficient.

\section{Implementation}

\begin{lstlisting}
#include <iostream>
#include <vector>
#include <string>
#include <climits>
using namespace std;

int main() {
    int charW, charH, numChars;
    cin >> charW >> charH >> numChars;

    vector<char> charLabel(numChars);
    vector<vector<string>> patterns(numChars,
        vector<string>(charH));

    for (int k = 0; k < numChars; k++) {
        cin >> charLabel[k];
        for (int r = 0; r < charH; r++)
            cin >> patterns[k][r];
    }

    int imgH, imgW;
    cin >> imgH >> imgW;
    vector<string> image(imgH);
    for (int r = 0; r < imgH; r++)
        cin >> image[r];

    int textRows = imgH / charH;
    int textCols = imgW / charW;

    for (int tr = 0; tr < textRows; tr++) {
        for (int tc = 0; tc < textCols; tc++) {
            int bestDist = INT_MAX;
            char bestChar = '?';

            for (int k = 0; k < numChars; k++) {
                int dist = 0;
                for (int r = 0; r < charH && dist < bestDist; r++) {
                    for (int c = 0; c < charW; c++) {
                        int ir = tr * charH + r;
                        int ic = tc * charW + c;
                        if (image[ir][ic] != patterns[k][r][c])
                            dist++;
                    }
                }
                if (dist < bestDist) {
                    bestDist = dist;
                    bestChar = charLabel[k];
                    if (dist == 0) break; // exact match
                }
            }
            cout << bestChar;
        }
        cout << "\n";
    }

    return 0;
}
\end{lstlisting}

\section{Example}

Suppose characters are $3 \times 5$ pixels (width $\times$ height), with
patterns for \texttt{A} and \texttt{B}:

\begin{verbatim}
A:          B:
.#.         ##.
#.#         #.#
###         ##.
#.#         #.#
#.#         ##.
\end{verbatim}

An input image of height 5, width 6 containing ``AB'':
\begin{verbatim}
.#.##.
#.##.#
####.
#.##.#
#.###.
\end{verbatim}

The decoder extracts each $3 \times 5$ block, matches against the patterns,
and outputs ``AB''.

\section{Notes}

\begin{itemize}
  \item The inner loop includes an early termination optimization: if the
    running distance already exceeds the best so far, we skip the remaining
    pixels for that pattern.
  \item When an exact match is found ($d = 0$), we immediately stop
    comparing further patterns.
  \item For larger alphabets, bitmask representations of pixel rows enable
    faster comparisons using XOR and popcount operations.
\end{itemize}

\end{document}
